\documentclass{article}
\usepackage{amsmath, amssymb, amsthm}
\title{Proof of Smoothness and Uniqueness for Morphic-Regularized Navier-Stokes Equations}
\author{Groby}
\date{\today}

\begin{document}

\maketitle

\section*{Introduction}

In this final step, we address the uniqueness of solutions to the morphic-regularized Navier-Stokes equations. While the introduction of the regularization term \( \epsilon \mathbf{R}(\mathbf{u}) \) provides smooth solutions, we now focus on proving that these solutions are unique, ensuring that the system yields a single, well-defined solution for any given set of initial conditions.

\section*{Step 5: Uniqueness of Solutions}

The uniqueness of the Navier-Stokes solutions is a central challenge. To prove uniqueness for the morphic-regularized system, we must show that if two solutions exist, their difference diminishes over time, ultimately converging to zero. This step involves comparing two potential solutions \( \mathbf{u}_1 \) and \( \mathbf{u}_2 \), and demonstrating that the difference between them vanishes.

\subsection*{Consider Two Solutions}

Let \( \mathbf{u}_1(x,t) \) and \( \mathbf{u}_2(x,t) \) be two smooth solutions to the morphic-regularized Navier-Stokes equations. We define their difference as:
\[
\mathbf{w}(x,t) = \mathbf{u}_1(x,t) - \mathbf{u}_2(x,t).
\]
Taking the difference between the two versions of the Navier-Stokes equations gives:
\[
\frac{\partial \mathbf{w}}{\partial t} + (\mathbf{u}_1 \cdot \nabla) \mathbf{w} - (\mathbf{u}_2 \cdot \nabla) \mathbf{w} = \nu \nabla^2 \mathbf{w} + \epsilon \mathbf{R}(\mathbf{w}).
\]
Since \( \mathbf{u}_1 \) and \( \mathbf{u}_2 \) are both incompressible, we also have:
\[
\nabla \cdot \mathbf{w} = 0.
\]

\subsection*{Energy Estimate for \( \mathbf{w}(x,t) \)}

We now derive an energy estimate for \( \mathbf{w}(x,t) \). Taking the \( L^2 \)-norm of the difference \( \mathbf{w} \) yields the energy:
\[
E_w(t) = \frac{1}{2} \int_{\mathbb{R}^n} |\mathbf{w}(x,t)|^2 \, dx.
\]
Taking the time derivative of \( E_w(t) \) gives:
\[
\frac{d}{dt} E_w(t) = \int_{\mathbb{R}^n} \mathbf{w}(x,t) \cdot \frac{\partial \mathbf{w}}{\partial t} \, dx.
\]
Substituting the equation for \( \frac{\partial \mathbf{w}}{\partial t} \), we get:
\[
\frac{d}{dt} E_w(t) = \int_{\mathbb{R}^n} \mathbf{w} \cdot \left( - (\mathbf{u}_1 \cdot \nabla) \mathbf{w} + (\mathbf{u}_2 \cdot \nabla) \mathbf{w} + \nu \nabla^2 \mathbf{w} + \epsilon \mathbf{R}(\mathbf{w}) \right) \, dx.
\]
The convective terms cancel out due to incompressibility, and the remaining terms yield:
\[
\frac{d}{dt} E_w(t) = - \nu \|\nabla \mathbf{w}\|_{L^2}^2 + \epsilon \int_{\mathbb{R}^n} \mathbf{w} \cdot \mathbf{R}(\mathbf{w}) \, dx.
\]

\subsection*{Dissipation of Energy Difference}

The viscous term \( - \nu \|\nabla \mathbf{w}\|_{L^2}^2 \) ensures dissipation of the energy difference over time. The regularization term \( \epsilon \mathbf{R}(\mathbf{w}) \) provides additional smoothing, dissipating any remaining energy in \( \mathbf{w} \). Since the regularization term is dissipative, we have:
\[
\int_{\mathbb{R}^n} \mathbf{w} \cdot \mathbf{R}(\mathbf{w}) \, dx \leq 0,
\]
which guarantees further dissipation.

Thus, the total dissipation is given by:
\[
\frac{d}{dt} E_w(t) \leq - \nu \|\nabla \mathbf{w}\|_{L^2}^2.
\]
This implies that the energy difference \( E_w(t) \) decreases over time, driving \( \mathbf{w}(x,t) \) toward zero.

\subsection*{Uniqueness of Solutions}

Since \( E_w(t) \geq 0 \) and \( \frac{d}{dt} E_w(t) \leq 0 \), it follows that \( E_w(t) \to 0 \) as \( t \to \infty \). Therefore, the difference between the two solutions \( \mathbf{u}_1 \) and \( \mathbf{u}_2 \) vanishes over time, meaning that:
\[
\mathbf{u}_1(x,t) = \mathbf{u}_2(x,t) \quad \text{for all} \quad t.
\]
This establishes that the morphic-regularized Navier-Stokes equations have a unique solution for all initial conditions.

\section{Conclusion of Step 5}

We have demonstrated that the morphic-regularized Navier-Stokes equations admit a unique smooth solution for all time. The dissipation introduced by the regularization term ensures that any two solutions will converge to the same result, proving uniqueness.

\end{document}
